\medskip

\emph{Pour chacune des affirmations suivantes, indiquer si elle est juste ou fausse.\\ Chaque réponse doit être justifiée.\\ Une réponse non justifiée ne rapporte aucun point.}

\medskip

\begin{enumerate}
\item Dans une classe de 24 élèves, il y a 14 filles et 10 garçons.

\textbf{Affirmation 1 :}

Il est possible de constituer 272 groupes différents de quatre élèves composés de deux filles et deux garçons.
\item Soit $f$ la fonction définie sur $\R$ par $f(x) = 3\sin(2x + \pi)$ et $C$ sa courbe représentative dans un repère donné.

\textbf{Affirmation 2 :}

Une équation de la tangente à $C$ au point d’abscisse $\dfrac{\pi}{2}$ est $y = 6x - 3\pi$.

\item On considère la fonction $F$ définie sur $]0~;~+\infty[$ par $F(x) = (2x + 1)\ln (x)$.

\textbf{Affirmation 3 :}

La fonction $F$ est une primitive de la fonction $f$ définie sur $]0~;~+\infty[$ par $f(x) = \dfrac2x$.
\item On considère la fonction $g$ définie sur $\R$ par $g(t) = 45\e^{0,06t} + 20$.

\textbf{Affirmation 4 :}

La fonction $g$ est l’unique solution de l'équation différentielle

$(E_1)\: ∶ y' + 0,06y = 1,2$ vérifiant $g(0) = 65$.

\item On considère l’équation différentielle :

\[(E_2) :\quad y' - y = 3\e^{0,4x}\]

où $y$ est une fonction positive de la variable réelle $x$, définie et dérivable sur $\R$ et $y'$ la fonction dérivée de la fonction $y$.

\textbf{Affirmation 5 :}

Les solutions de l’équation $(E_2)$ sont des fonctions convexes sur $\R$.
\end{enumerate}

\medskip
